\subsection{Transferência de órbita}
\label{sec:hohmann}
A manobra de dois impulsos, denominada transferência de Hohmann, no contexto de RVD/B, é utilizada para ajustar a altitude orbital do \emph{chaser} em relação à órbita do \emph{target}.
Inicialmente, a espaçonave visitante parte de uma órbita circular de raio $r_1$ até uma órbita circular final de raio $r_2$. Além disso, os impulsos são instantâneos para minimizar o gasto de $\Delta v$ para a transferência.

\begin{figure}[ht]
\centering
\includegraphics[width=0.8\linewidth]{2 Formulacao do problema/Imagens/TransfHohmann.png}
\caption{Transferência de Hohmann. Fonte: Autor, adaptado de MATLAB.}
    \label{fig:transf_hohmann}
\end{figure}

\subsubsection{Descrição da manobra}
Em determinado ponto do perigeu, aplica-se o primeiro impulso, denominado $\Delta v_1$.
A espaçonave irá sair da órbita circular de raio $r_1$ e entrará em uma órbita elíptica de transferência, cujo perigeu está em $r_1$ e o apogeu, em $r_2$. 
% No caso da órbita circular, o perigeu e o apogeu são iguais.
No segundo impulso ($\Delta v_2$), aplicado no apogeu, a órbita é circularizada em $r_2$.
A figura \eqref{fig:transf_hohmann} ilustra o processo.

\subsubsection{Velocidades orbitais}
A velocidade inicial\fn{Apresentadas em $[km/s]$} na órbita circular de raio $r_1$ é similar ao que foi apresentado em \eqref{eq:veloc_insercao_orbital}:
\begin{equation}
    v_1 = \sqrt{\frac{\mu}{r_1}}.
\end{equation}

A velocidade no perigeu da órbita de transferência, de forma similar a \eqref{eq:vis_viva}, é dada por:

\begin{equation}
v_{Transf(1)}
=
\sqrt{2\mu \left(\frac{1}{r_1}
-\frac{1}{2a_{Transf}}\right)},
\end{equation}

onde o semi-eixo maior da órbita de transferência é encontrado conforme:
\begin{equation}
a_{Transf}
=
\frac{r_1 + r_2}{2}
\end{equation}

A variação de velocidade requerida no primeiro impulso é
\begin{equation}
\Delta v_1
=
v_{Transf(1)}
-v_1.
\end{equation}

De forma análoga, a velocidade na órbita circular de raio $r_2$ é
\begin{equation}
v_2
=
\sqrt{\frac{\mu}{r_2}},
\end{equation}

enquanto que a velocidade no apogeu da órbita de transferência é:

\begin{equation}
v_{Transf(2)}
=
\sqrt{2\mu \left(\frac{1}{r_2}
-\frac{1}{2a_{Transf}}\right)}.
\end{equation}

O segundo impulso é então
\begin{equation}
\Delta v_2
=
v_2
-v_{Transf(2)}.
\end{equation}

\subsubsection{Tempo de transferência e custo total}
O tempo total, da transferência orbital corresponde a metade do período da órbita elíptica:

\begin{equation}
t_{Transf}
=
\pi \sqrt{\frac{a_{Transf}^3}{\mu}}.
\end{equation}

Para encontrar o custo de velocidade total da manobra, utiliza-se:

\begin{equation}
    \Delta v_{Hohmann} = \Delta v_1 + \Delta v_2.
\end{equation}

% \subsubsection{Considerações finais}
% Após o término da transferência de Hohmann, o \emph{chaser} encontra-se em uma órbita coplanar e próxima da órbita do \emph{target}, no entanto, estará defasado em fase orbital.
% Se necessário, ajustes podem ser realizados para sincronizar ou reduzir as diferenças entre as anomalias verdadeiras ($\nu$) das espaçonaves.
% Após isso, inicia-se o cálculo da dinâmica do movimento relativo, descrita pelas Equações de Hill/Clohessy--Wiltshire.